\documentclass[a4paper,12pt, listof=nochaptergap, oneside]{book}
\usepackage{program}

\begin{document}
\frontmatter
	\include{Depan/1.Sampul}
	\include{Depan/2.Judul}
	\include{Depan/3.Persetujuan}
	\include{Depan/4.Identitas}
	\include{Depan/5.Orisinilitas}
	\include{Depan/6.Riwayat}
	\include{Depan/7.Ringkasan}
	\include{Depan/8.Summary}
	\include{Depan/9.Pengantar}
	\include{Depan/10.DaftarIsi}
	\include{Depan/11.DaftarTabel}
	\include{Depan/12.DaftarGambar}
	\include{Depan/13.DaftarSimbol}
	\include{Depan/14.DaftarSingkatan}
	\include{Depan/15.DaftarLampiran}
%Penomoran Halaman
\fancyhf{}
\cfoot{\thepage}
\renewcommand{\headrulewidth}{0pt}
%\pagestyle{fancy}
%\fancypagestyle{plain}
%{\fancyhf{}
%	\renewcommand{\headrulewidth}{0pt}
%	\cfoot{\thepage}}

%Memunculkan BAB
	\makeatletter
	\renewcommand{\@makechapterhead}[1]{
		\vspace*{-15\p@}%Mengatur jarak bab
		
		{\parindent \z@ \raggedright \ifnum \c@secnumdepth >\m@ne
			\if@mainmatter
			\center \MakeUppercase{\large\bfseries \@chapapp\space \thechapter}
			%\chapapp\space memunculkan kata "Bab"
			\par\nobreak
			\vskip 1\p@ % Mengatur jarak antara "Bab 1" dan "Pendahuluan"
			\fi
			\fi
			\interlinepenalty\@M
			\large \uppercase{\bfseries #1\par\nobreak} %Bagian ini memunculkan judul bab
			\vskip 24\p@ % Mengatur jarak antara "Bab 1" dengan "Teks"
	}}
	\makeatother
%%%%%%%%%%%%%%%%%%%%%%%%%%%%%%%%%%%%%%%%%%%%%%%%%%%%%%%%%%%%%%%%%%%%%%%%%%%%%%%%%%
%\makeatletter
%\pretocmd{\chapter}{\addtocontents{toc}{\protect\addvspace{15\p@}}}{}{}
%\makeatother

\mainmatter	
\pagestyle{plain}
	\include{Isi/1.Pendahuluan}
	\include{Isi/2.Tinjauan}
	\include{Isi/3.Kerangka}
	\include{Isi/4.Metode}
	\include{Isi/5.Hasil}
	\include{Isi/6.Penutup}
%%%%%%%%%%%%%%%%%%%%%%%%%%%%%%%%%%%%%%%%%%%%%%%%%%%%%%%%%%%%%%%%%%%%%%%%%%%%%%%%%%%	

\backmatter
	\include{Belakang/1.Pustaka}
	\include{Belakang/2.Lampiran}
\end{document}
